%% kapa: exam=1 round=21 year=2010 period=2 form=A subject=accounting subject_ko=회계학 seq=5 number=5
%% source: Q-Net 감정평가사 기출문제 (cst00309) artlSeq=5024181 file=extracted/1008/5024181/2010년 제21회 감정평가사 1차시험 기출문제/2교시.hwp
%% conversion: kordoc 4.24.51 -> tools/gen1.py (auto); review=proofread-sonnet; flags=none
\begin{question}{5}
(주)한국은 재고자산 수량결정과 관련하여 계속기록법을 채택하고 있다. 다음은 계속기록법으로 작성한 (주)한국의 20×0년의 매출원가와 관련된 자료이다.

{\small\centering
\begin{tabular}{|l|r|r|r|}
\hline
 & \multicolumn{1}{c|}{수량} & \multicolumn{1}{c|}{단가} & \multicolumn{1}{c|}{합계} \\
\hline
20×0년 초 재고자산 & 100개 & ￦300 & ￦30,000 \\
\hline
20×0년 매입액 & 200개 & 300 & 60,000 \\
\hline
20×0년 말 재고자산 & 150개 & 300 & 45,000 \\
\hline
\end{tabular}\par}

(주)한국이 20×0년 말 재고자산을 실사한 결과 재고자산감모수량이 30개(재고자산감모손실 ￦9,000) 발생하였다. 또한, (주)한국은 20×0년 말 재고자산의 단위당 순실현가능가치가 ￦200으로 하락하여 재고자산평가손실을 인식하여야 한다. (주)한국이 20×0년도에 인식할 재고자산평가손실은 얼마인가? (단, 재고자산감모손실을 먼저 인식한 후 재고자산평가손실을 산출함을 가정하고, 기초재고자산에 대한 평가충당금은 없었음)
\begin{choices}
  \item ￦6,000
  \item ￦9,000
  \item ￦12,000
  \item ￦15,000
  \item ￦18,000
\end{choices}
\end{question}
\answer{3}
